\section{Distributed Training: las Rare Failure se vuelven eventos cotidianos}

Un Frontier run depende al mismo tiempo de gran cantidad de accelerators, host, collective network, dataset storage, checkpoint storage, scheduler y cloud quota. Aunque la tasa de fallas de cada componente sea muy baja, al aumentar la escala las interrupciones se vuelven frecuentes. El objetivo del sistema no es aspirar ilusoriamente a cero fallas, sino acotar el failure domain, detectar rápido, conservar el estado científico y recuperar el avance efectivo.

\subsection{Workers, network, storage y scheduler}

La \term{collective communication} (comunicación colectiva) abarca las primitive de coordinación entre múltiples dispositivos, como all-reduce, all-gather y reduce-scatter; un solo worker lento o una network partition basta para frenar el entrenamiento síncrono. El storage debe suministrar token de forma continua y además recibir checkpoint de gran tamaño; el scheduler, por su parte, mantiene la membership, la placement y el restart. Toda dependencia necesita health signal, timeout, backpressure y ownership.

\lecturefigure{05-training-failure-domains.png}{El Frontier training acopla cuatro tipos de failure domain: workers, network, storage y scheduler.}{Subtítulos locales 15:00--17:20; redibujo conceptual.}

\begin{knowledgebox}{Lectura de la figura: primero distinguir hard failure de silent corruption}
Las hard failure incluyen la terminación de procesos, la desconexión de dispositivos y el storage timeout, y suelen activar un restart; la silent corruption puede manifestarse como datos erróneos, collective bit error, stale code o gradientes anómalos, y es más peligrosa. El monitoreo no solo debe responder «¿el job sigue en ejecución?», sino también «¿sigue ejecutando el mismo experimento válido?».
\end{knowledgebox}

\begin{importantbox}{Useful training rate}
Sea $F_{peak}$ el compute pico teórico, $u$ la tasa de utilización efectiva y $d$ la proporción perdida por fallas, recuperación e inactividad; entonces la useful rate a largo plazo puede aproximarse como
\[
F_{useful}=F_{peak}\,u(1-d).
\]
Aumentar el kernel MFU solo mejora $u$; si el checkpoint es demasiado lento, la detección de fallas es tardía o la recuperación es inconsistente y por ello $d$ resulta muy grande, el hardware costoso sigue sin traducirse en avance efectivo del entrenamiento.
\end{importantbox}

\subsection{Resumen de la sección}

La escala vuelve habituales las rare failure. El Distributed training debe optimizar a la vez la steady-state utilization y el recovery tax, y usar la consistencia del experimento, y no el «job resumed», como criterio de recuperación exitosa.
